\subsection{An exact next-step criterion with frozen ordinary support}

\begin{result}{E.146}{Proved}{The quotient gate is sufficient for one finite step}
Use the groups \(Q_{n-1},Q_n,B_n\) of E.144, with
\(n\ge2\). Let
\[
 S=\{3,\infty,q_1,\ldots,q_t\},\qquad
 q_i\equiv-1\pmod3,
\]
and include five and eleven among the \(q_i\).
Fix a lower homomorphism \(\rho\) and actual ordinary
point classes at all the old places, as in E.144.
Assume its inertias at five and eleven are,
respectively, nonzero multiples of the two marked
generators in the Frattini quotient \(F/P_2\).
Keep the actual arithmetic projection fixed.

Then the following conditions are equivalent:
\begin{enumerate}[label=(\roman*)]
\item There is a lift to \((Q_n/B_n)\rtimes D\),
unramified outside \(S\), matching all the specified
old point classes.
\item There is a lift to \(Q_n\rtimes D\), unramified
outside \(S\) and a finite set of new primes, which
matches those same old point classes and is the
image of an ordinary point at every new prime.
\end{enumerate}
When (i) holds, at most two new primes are needed.
They can be chosen outside any prescribed finite
set and satisfy
\[
                         q\equiv2\pmod9.
 \tag{E.146.1}
\]
At each new prime the point may be taken to be
\(q^m\) or \(1-q^m\), with
\[
              m\in\{3^{\,n-1},\,2\cdot3^{\,n-1}\}.
 \tag{E.146.2}
\]
Arbitrarily many additional pairs of such new primes
may also be introduced without changing any old
point class.

Thus the actual ordinary integer valuations at old
support remain frozen as externally specified point
data, and every new valuation can
be fixed as a positive ordinary integer when it is
introduced. The finite quotient itself detects those
valuations only modulo \(3^n\).
Matching the specified actual point at three is an
additional requirement of this construction; it is
not forced by boundary rigidity A.31 or by ordinary
pro-three abelianization.
The statement is one finite step. It
does not assert that a choice satisfying the next
quotient gate can always be retained at every later
level, or that the resulting finite abelianization
is the reduction of one completed pair ordinary at
every zero-support place.

There is also a version with no point condition at
three. In that version the quotient gate remains
equivalent to the next-step condition, and no new
prime is needed to lift while matching the old
support. Moreover the minus-character part of
the quotient matching gate is automatically
soluble. Only its plus-character Frobenius
conditions remain: in each plus coordinate
the old local discrepancies must be restrictions
of one global cubic character. New support
can still be introduced if desired.
\end{result}
\begin{proof}
Necessity is E.144, including its requirement of
actual local point images. We prove sufficiency.
Write \(\bar Q=Q_n/B_n\), and fix a lift
\(\bar\rho\) as in (i), together with a choice of
each old local aligning frame.

Enlarge the finite arithmetic quotient \(D\), if
necessary, to retain \(\mu_9\) and the actual
finite cusp-zero and cusp-one path values in
\(Q_n\). This is an enlargement of the fixed
arithmetic data, not a modification of \(\rho\):
adjoin the actual arithmetic representation to
the given homomorphism. These additional data
are unramified outside three by the actual
based-action and tangential-path statements
in E.126 and E.128. Denote the enlarged
arithmetic quotient by \(D'\).

First lift \(\bar\rho\) over the old ramification
set, without yet imposing all old point matches.
The kernel \(B_n\) is
\(\mathbb F_3(1)^2\); its action is independent
of the geometric part of \(\bar\rho\), since
it is geometrically central. The extension
obstruction belongs to
\[
 H^2(G_{\mathbb Q,S},\mu_3)^2
             =\operatorname{Br}(\mathbb Z[S^{-1}])[3]^2.
 \tag{E.146.3}
\]
The equality uses the same finite comparison and
the vanishing Picard group as E.144.
Every old local restriction is zero because the
specified actual point supplies a lift.
The Brauer local invariant map is injective on
(E.146.3). Hence there is a global weak lift
\(t:G_{\mathbb Q,S}\to Q_n\rtimes D'\).

With the old aligning frames fixed, the discrepancy
between \(t\) and each prescribed local point is
a well-defined class of \(H^1(\mathbb Q_v,B_n)\).
Changing the lift of an aligning frame changes
only a kernel coboundary. For \(q_i\ne3\),
Kummer theory identifies
\[
 H^1(\mathbb Q_{q_i},\mu_3)
    =\mathbb Q_{q_i}^{\times}/
                   (\mathbb Q_{q_i}^{\times})^3
    =\mathbb F_3\cdot\kappa(q_i).
 \tag{E.146.4}
\]
Indeed every unit is a cube because
\(q_i-1\) is prime to three.
Twist \(t\), in each of the two \(B_n\) coordinates,
by suitable global combinations of the
\(\kappa(q_i)\). This attains every old finite
point match away from three.
Each \(\kappa(q_i)\) restricts to zero at the
other old primes different from three.
Then use the two \(\kappa(3)\) coefficients to
correct the valuation coordinates at three.
This leaves, at three only, two unit discrepancies.
The real first-cohomology group is zero.

The unit quotient of \(\mathbb Q_3^\times\)
modulo cubes has order three and is generated
by \(\kappa(2)\). Thus the remaining corrections
are \(c_x\kappa(2)\) and \(c_y\kappa(2)\), with
\(c_x,c_y\in\mathbb F_3\).
For every prime satisfying (E.146.1),
\[
               \kappa(q)|_{\mathbb Q_3}=\kappa(2),
 \tag{E.146.5}
\]
because \(q/2\in1+9\mathbb Z_3\) is a cube.
Also \(v_3(q^2-1)=1\).

We next show that primes satisfying (E.146.1)
can simultaneously carry the required cusp
Frobenius for \(\bar\rho\).
On the kernel of the arithmetic quotient \(D'\),
the geometric coordinate of \(\bar\rho\) is
a homomorphism to \(\bar Q\).
Its image contains the two independent Frattini
inertia vectors at five and eleven, since
\(D'\) is unramified there.
The finite three-group Frattini argument gives
the entire \(\bar Q\) as its image.
Consequently, above every element \(d\in D'\),
the joint image contains every element of
the geometric fibre.
Choose \(d\) with cyclotomic value two modulo
nine. The actual cyclotomic representation is
surjective onto \((\mathbb Z/9)^\times\), so
such a \(d\) exists.
In its geometric fibre choose the cusp-zero
value or the cusp-one value, as required.
Chebotarev applied to this finite joint field
produces arbitrarily large primes with that
Frobenius and with (E.146.1)
\citep[Theorem 8.31]{CohMilneANT}.
In particular, all previously specified primes
can be avoided.

If \(c_x\ne0\), choose one such prime \(q_x\)
with cusp-zero Frobenius and twist the lift
by \(c_x\kappa(q_x)\) in the
\(x^{3^{n-1}}\) coordinate of \(B_n\).
If \(c_y\ne0\), choose a distinct such prime
\(q_y\) with cusp-one Frobenius and twist
by \(c_y\kappa(q_y)\) in the other coordinate.
By (E.146.5) these correct the two remaining
unit classes at three.
At every other old finite place these twists
are unit Kummer classes and hence zero by
(E.146.4). Thus all old point matches remain
valid.

It remains to prove the complete new local
point assertion, including Frobenius.
At a new prime \(q\), the chosen quotient
\(\bar\rho\) is unramified and has the selected
cusp Frobenius. The lift before the new twists
is also unramified at \(q\).
The unramified first-cohomology group of the
kernel is
\[
 H^1_{\mathrm{nr}}(\mathbb Q_q,B_n)
     = B_n/(q-1)B_n=0.
 \tag{E.146.6}
\]
Therefore all unramified lifts of this aligned
quotient cusp homomorphism are conjugate by
the kernel. The full local first-cohomology
group of \(B_n\) is determined by inertia,
as in (E.146.4).
After the twist, inertia is exactly
\[
 x^{c_x3^{n-1}}\quad\hbox{or}\quad
 y^{c_y3^{n-1}}.
 \tag{E.146.7}
\]
Twists from the other new prime are locally
unit Kummer classes and do not change this
local cohomology class.

For clarity, the actual point class having
(E.146.7) is available before taking the finite
quotient. Twist the appropriate local cusp
section by the full local Kummer cocycle
of \(q^m\), with \(m\) as in (E.146.2).
Its pro-three abelianization agrees with the
point \(q^m\), or \(1-q^m\), because the
complementary unit has trivial pro-three
Kummer class at \(q\equiv-1\pmod3\).
The boundary-fibre rigidity A.31 identifies
this section with that actual point class.
Its image in \(\bar Q\) is the selected cusp
section, since its peripheral powers lie
in \(B_n\).
Its kernel discrepancy is the inertia class
in (E.146.7).
Equations (E.146.4) and (E.146.6) therefore
identify the entire new local homomorphism,
including Frobenius, with the actual point
image. No arbitrary unramified Frobenius
is being discarded.

The lift is unramified outside \(S\) and these
at most two primes: the initial lift was
unramified outside \(S\), and each of the
Kummer twists is ramified only at its prime
and possibly at three.
It projects to the fixed literal lower
homomorphism throughout.
All changes to local frames have occurred
inside the indicated kernel and are allowed
geometric conjugacies.

Finally, extra new support can be forced
without changing old conditions.
Choose two further primes, both with the
same cusp Frobenius and both two modulo nine.
Use coefficients one and two in the same
\(B_n\) coordinate.
Their local unit corrections at three sum
to zero, and their restrictions at every
other old place are zero.
At the two new primes the preceding argument
gives ordinary valuations \(3^{n-1}\) and
\(2\cdot3^{n-1}\), both positive integers.
Repeating with disjoint pairs proves the
last assertion. This construction establishes
the exact one-step equivalence, not the
nonemptiness of its quotient gate at every
subsequent level.

Finally omit the point requirement at three.
For either the minus-character part of the
quotient obstruction or the subsequent
\(B_n\)-obstruction, its Brauer invariants
at every old prime different from three
are zero, because those prescribed points
lift. All other invariants except possibly
three are zero by unramifiedness and the
absence of odd-primary real invariants.
Their sum is zero, so the invariant at
three also vanishes. Thus both obstructions
vanish over the old \(S\).
For every minus coordinate the old local
first-cohomology discrepancies are valuation
classes; the independent global
\(\kappa(q_i)\) attain all of them.
There is no remaining local-three condition.
This proves the claimed automatic minus
matching and supplies the \(B_n\)-lift with
no extra primes. For a plus coordinate,
E.144 gives the single global parameter
\(\bar\eta_4\). Its local first cohomology
at each \(q_i\equiv-1\pmod3\) is unramified,
so its full remaining condition is exactly
the stated Frobenius localization condition.
One can add fresh cusp-selected primes by
the same kernel twists; their correction
at three need no longer be cancelled.
\end{proof}
